\documentclass[../main.tex]{subfiles}
\begin{document}

\chapter{データセンタ技術}
\lessoninfo{
  video={P4L4},
  textbook={Silberschatz ら \cite{silberschatz2018} ch.~1；OSTEP \cite{ostep} ch.~48；Armbrust ら \cite{armbrust2009}；Mell と Grance \cite{mell2011}},
  keywords={インターネットサービス，多層アーキテクチャ，スケールアウト，ロードバランサ，ホモジニアス構成，ヘテロジニアス構成，クラウドコンピューティング，伸縮性，大数の法則，ユーティリティコンピューティング，IaaS，PaaS，SaaS，マルチテナンシ},
}

これまでの章は，異なるマシン上のプロセスが協調するための機構---遠隔手続き
呼び出し，分散ファイルシステム，そして第15章の分散共有メモリ---を扱って
きた．これらの機構は，より前の章のスレッド，サーバ，スケジューリング，
仮想化とともに，インターネット越しに使われるサービスを動かすデータセンタに
おいて，最大の規模で用いられている．この最後の章では，データセンタ技術を
大まかに概観する．インターネットサービスがどう構成され，多数のマシンへ
スケールアウトされるか，クラウドコンピューティングとは何か，それがなぜ
成り立ち，どのようなモデルを提供するか，基盤となるシステムに何を要求するか，
そしてどのような技術がそれを可能にするかである．

\section{インターネットサービス}
\label{sec:l16-services}

日常的に使われるアプリケーションの多くは，その利用者のマシン上では動作
しない．\source{P4L4，レッスンの概要，インターネットサービス．}それらは
\term[でーたせんた]{データセンタ}[datacenter]のマシン上で動作する．データ
センタとは，ネットワークで結ばれた多数の計算機をそのストレージ，電源，冷却
設備とともに収容し，計算サービスを提供するために運用する施設である．利用者は
インターネット越しにこれらのアプリケーションに到達する．

\begin{definition}[インターネットサービス]
  \term[いんたーねっとさーびす]{インターネットサービス}[Internet service]
  とは，Web インタフェースを通じてインターネット越しに利用者へ提供される
  あらゆるサービスである．
\end{definition}

Web メール，オンラインショップ，検索エンジン，ソーシャルネットワークは
インターネットサービスである．ほとんどのインターネットサービスは3つの構成
要素からなる（\cref{fig:l16-tiers}）．
\begin{description}
  \item[プレゼンテーション] 利用者とやり取りする．その要求を受け取り，Web
    ページを返し，ページのレイアウトや画像のように要求によらず同じである
    \emph{静的コンテンツ}を提供する．
  \item[ビジネスロジック] サービスの機能を実装する．各要求を処理し，検索の
    結果やショッピングカートの中身のような\emph{動的コンテンツ}を生成する．
  \item[データベース] 利用者，商品，注文といったサービスのデータを格納し，
    ビジネスロジックのためにそれらを取り出す．
\end{description}

\begin{definition}[多層アーキテクチャ]
  一連の構成要素として構成され，各構成要素が次の構成要素のサービスを利用
  するサービスは，\term[たそうあーきてくちゃ]{多層アーキテクチャ}%
  [multi-tier architecture]を持つ．各構成要素が1つの\emph{層}であり，上記の
  3つの構成要素は三層アーキテクチャを成す．
\end{definition}

層は論理的な構成要素である．別々のプロセスである必要はなく，別々のプロセス
が別々のマシンで動作する必要もない．Web サーバ，アプリケーションサーバ，
データベースサーバといった多数のオープンソースおよび商用のミドルウェア製品
が個々の層と層どうしのやり取りを実装しており，さまざまな構成で組み合わせ
られる．層が別々のプロセスとして動作する場合，プロセスは何らかの
IPC で通信する．同じマシン上ではメッセージパッシングや共有メモリ（第9章），
マシンをまたぐときは RPC や RMI（第13章）である．

\begin{figure}[htbp]
  \centering
  % Three-tier Internet service: (a) all tiers on one machine, (b) each tier
% on its own machine, the tiers communicating by RPC/RMI.
% Usage: % Three-tier Internet service: (a) all tiers on one machine, (b) each tier
% on its own machine, the tiers communicating by RPC/RMI.
% Usage: % Three-tier Internet service: (a) all tiers on one machine, (b) each tier
% on its own machine, the tiers communicating by RPC/RMI.
% Usage: \input{figures/l16-tiers}   (width ~118 mm)
\begin{tikzpicture}[x=1mm, y=1mm, font=\scriptsize\sffamily,
  >={Stealth[length=1.5mm]},
  tier/.style={draw, fill=ln-box, minimum width=22mm, minimum height=10mm, inner sep=1pt, align=center},
  cl/.style={draw, fill=white, minimum width=14mm, minimum height=10mm, inner sep=1pt, align=center},
  machine/.style={draw, rounded corners=0.8mm},
  lab/.style={font=\scriptsize\sffamily, text=ln-muted, align=center},
  who/.style={font=\scriptsize\sffamily, text=ln-muted, anchor=north west, inner sep=0.8pt},
  ttl/.style={font=\footnotesize\sffamily\bfseries, anchor=west}]
  % ---------------- (a) one machine
  \node[ttl] at (0,57) {\iflangja{(a) すべての層が1台のマシン上にある}{(a) all tiers on one machine}};
  \node[cl] (c1) at (7,40) {\iflangja{クライアント}{clients}};
  \draw[machine] (20,30) rectangle (118,53);
  \node[who] at (20.5,52.5) {\iflangja{マシン}{machine}};
  \node[tier] (p1) at (33,42) {\iflangja{プレゼンテーション}{presentation}\\[0.3mm]\textcolor{ln-muted}{\iflangja{静的コンテンツ}{static content}}};
  \node[tier] (b1) at (69,42) {\iflangja{ビジネスロジック}{business logic}\\[0.3mm]\textcolor{ln-muted}{\iflangja{動的コンテンツ}{dynamic content}}};
  \node[tier] (d1) at (105,42) {\iflangja{データベース}{database}\\[0.3mm]\textcolor{ln-muted}{\iflangja{データストア}{data store}}};
  \draw[<->, thick] (c1.east |- p1) -- (p1.west);
  \draw[<->, thick] (p1) -- (b1);
  \draw[<->, thick] (b1) -- (d1);
  \node[lab] at (69,33.2)
    {\iflangja{1つのプロセス，または IPC・共有メモリで通信する複数のプロセス}{one process, or several processes communicating by IPC or shared memory}};
  % ---------------- (b) separate machines
  \node[ttl] at (0,24) {\iflangja{(b) 各層が別のマシン上にある}{(b) each tier on its own machine}};
  \node[cl] (c2) at (7,9) {\iflangja{クライアント}{clients}};
  \foreach \x/\n in {33/{\iflangja{Web サーバ}{web server}}, 69/{\iflangja{アプリケーションサーバ}{application server}}, 105/{\iflangja{データベースサーバ}{database server}}} {
    % Japanese: 2 mm wider boxes, or "アプリケーションサーバ" touches the frame.
    \draw[machine] (\x-\iflangja{14}{13},0) rectangle (\x+\iflangja{14}{13},20);
    \node[who] at (\x-12.5,19.5) {\n};
  }
  \node[tier] (p2) at (33,8) {\iflangja{プレゼンテーション}{presentation}};
  \node[tier] (b2) at (69,8) {\iflangja{ビジネスロジック}{business logic}};
  \node[tier] (d2) at (105,8) {\iflangja{データベース}{database}};
  \draw[<->, thick] (c2.east |- p2) -- (p2.west);
  \draw[<->, thick] (p2) -- (b2);
  \draw[<->, thick] (b2) -- (d2);
  \node[lab, anchor=south] at (51,8.6) {RPC};
  \node[lab, anchor=north] at (51,7.4) {RMI};
  \node[lab, anchor=south] at (87,8.6) {RPC};
  \node[lab, anchor=north] at (87,7.4) {RMI};
\end{tikzpicture}
   (width ~118 mm)
\begin{tikzpicture}[x=1mm, y=1mm, font=\scriptsize\sffamily,
  >={Stealth[length=1.5mm]},
  tier/.style={draw, fill=ln-box, minimum width=22mm, minimum height=10mm, inner sep=1pt, align=center},
  cl/.style={draw, fill=white, minimum width=14mm, minimum height=10mm, inner sep=1pt, align=center},
  machine/.style={draw, rounded corners=0.8mm},
  lab/.style={font=\scriptsize\sffamily, text=ln-muted, align=center},
  who/.style={font=\scriptsize\sffamily, text=ln-muted, anchor=north west, inner sep=0.8pt},
  ttl/.style={font=\footnotesize\sffamily\bfseries, anchor=west}]
  % ---------------- (a) one machine
  \node[ttl] at (0,57) {\iflangja{(a) すべての層が1台のマシン上にある}{(a) all tiers on one machine}};
  \node[cl] (c1) at (7,40) {\iflangja{クライアント}{clients}};
  \draw[machine] (20,30) rectangle (118,53);
  \node[who] at (20.5,52.5) {\iflangja{マシン}{machine}};
  \node[tier] (p1) at (33,42) {\iflangja{プレゼンテーション}{presentation}\\[0.3mm]\textcolor{ln-muted}{\iflangja{静的コンテンツ}{static content}}};
  \node[tier] (b1) at (69,42) {\iflangja{ビジネスロジック}{business logic}\\[0.3mm]\textcolor{ln-muted}{\iflangja{動的コンテンツ}{dynamic content}}};
  \node[tier] (d1) at (105,42) {\iflangja{データベース}{database}\\[0.3mm]\textcolor{ln-muted}{\iflangja{データストア}{data store}}};
  \draw[<->, thick] (c1.east |- p1) -- (p1.west);
  \draw[<->, thick] (p1) -- (b1);
  \draw[<->, thick] (b1) -- (d1);
  \node[lab] at (69,33.2)
    {\iflangja{1つのプロセス，または IPC・共有メモリで通信する複数のプロセス}{one process, or several processes communicating by IPC or shared memory}};
  % ---------------- (b) separate machines
  \node[ttl] at (0,24) {\iflangja{(b) 各層が別のマシン上にある}{(b) each tier on its own machine}};
  \node[cl] (c2) at (7,9) {\iflangja{クライアント}{clients}};
  \foreach \x/\n in {33/{\iflangja{Web サーバ}{web server}}, 69/{\iflangja{アプリケーションサーバ}{application server}}, 105/{\iflangja{データベースサーバ}{database server}}} {
    % Japanese: 2 mm wider boxes, or "アプリケーションサーバ" touches the frame.
    \draw[machine] (\x-\iflangja{14}{13},0) rectangle (\x+\iflangja{14}{13},20);
    \node[who] at (\x-12.5,19.5) {\n};
  }
  \node[tier] (p2) at (33,8) {\iflangja{プレゼンテーション}{presentation}};
  \node[tier] (b2) at (69,8) {\iflangja{ビジネスロジック}{business logic}};
  \node[tier] (d2) at (105,8) {\iflangja{データベース}{database}};
  \draw[<->, thick] (c2.east |- p2) -- (p2.west);
  \draw[<->, thick] (p2) -- (b2);
  \draw[<->, thick] (b2) -- (d2);
  \node[lab, anchor=south] at (51,8.6) {RPC};
  \node[lab, anchor=north] at (51,7.4) {RMI};
  \node[lab, anchor=south] at (87,8.6) {RPC};
  \node[lab, anchor=north] at (87,7.4) {RMI};
\end{tikzpicture}
   (width ~118 mm)
\begin{tikzpicture}[x=1mm, y=1mm, font=\scriptsize\sffamily,
  >={Stealth[length=1.5mm]},
  tier/.style={draw, fill=ln-box, minimum width=22mm, minimum height=10mm, inner sep=1pt, align=center},
  cl/.style={draw, fill=white, minimum width=14mm, minimum height=10mm, inner sep=1pt, align=center},
  machine/.style={draw, rounded corners=0.8mm},
  lab/.style={font=\scriptsize\sffamily, text=ln-muted, align=center},
  who/.style={font=\scriptsize\sffamily, text=ln-muted, anchor=north west, inner sep=0.8pt},
  ttl/.style={font=\footnotesize\sffamily\bfseries, anchor=west}]
  % ---------------- (a) one machine
  \node[ttl] at (0,57) {\iflangja{(a) すべての層が1台のマシン上にある}{(a) all tiers on one machine}};
  \node[cl] (c1) at (7,40) {\iflangja{クライアント}{clients}};
  \draw[machine] (20,30) rectangle (118,53);
  \node[who] at (20.5,52.5) {\iflangja{マシン}{machine}};
  \node[tier] (p1) at (33,42) {\iflangja{プレゼンテーション}{presentation}\\[0.3mm]\textcolor{ln-muted}{\iflangja{静的コンテンツ}{static content}}};
  \node[tier] (b1) at (69,42) {\iflangja{ビジネスロジック}{business logic}\\[0.3mm]\textcolor{ln-muted}{\iflangja{動的コンテンツ}{dynamic content}}};
  \node[tier] (d1) at (105,42) {\iflangja{データベース}{database}\\[0.3mm]\textcolor{ln-muted}{\iflangja{データストア}{data store}}};
  \draw[<->, thick] (c1.east |- p1) -- (p1.west);
  \draw[<->, thick] (p1) -- (b1);
  \draw[<->, thick] (b1) -- (d1);
  \node[lab] at (69,33.2)
    {\iflangja{1つのプロセス，または IPC・共有メモリで通信する複数のプロセス}{one process, or several processes communicating by IPC or shared memory}};
  % ---------------- (b) separate machines
  \node[ttl] at (0,24) {\iflangja{(b) 各層が別のマシン上にある}{(b) each tier on its own machine}};
  \node[cl] (c2) at (7,9) {\iflangja{クライアント}{clients}};
  \foreach \x/\n in {33/{\iflangja{Web サーバ}{web server}}, 69/{\iflangja{アプリケーションサーバ}{application server}}, 105/{\iflangja{データベースサーバ}{database server}}} {
    % Japanese: 2 mm wider boxes, or "アプリケーションサーバ" touches the frame.
    \draw[machine] (\x-\iflangja{14}{13},0) rectangle (\x+\iflangja{14}{13},20);
    \node[who] at (\x-12.5,19.5) {\n};
  }
  \node[tier] (p2) at (33,8) {\iflangja{プレゼンテーション}{presentation}};
  \node[tier] (b2) at (69,8) {\iflangja{ビジネスロジック}{business logic}};
  \node[tier] (d2) at (105,8) {\iflangja{データベース}{database}};
  \draw[<->, thick] (c2.east |- p2) -- (p2.west);
  \draw[<->, thick] (p2) -- (b2);
  \draw[<->, thick] (b2) -- (d2);
  \node[lab, anchor=south] at (51,8.6) {RPC};
  \node[lab, anchor=north] at (51,7.4) {RMI};
  \node[lab, anchor=south] at (87,8.6) {RPC};
  \node[lab, anchor=north] at (87,7.4) {RMI};
\end{tikzpicture}

  \caption{三層のインターネットサービス．(a) すべての層が1台のマシン上で，
    1つのプロセスとして，あるいは複数のプロセスとして動作する．(b) 各層が
    自分のマシンで動作し，層どうしは RPC または RMI で通信する．}
  \label{fig:l16-tiers}
\end{figure}

\section{スケールアウト構成}
\label{sec:l16-scaleout}

1台のマシンが毎秒処理できる要求の数は，サーバソフトウェアをどれほど効率
よく設計しても限られる（第6章）．\source{P4L4，インターネットサービスの
構成，ホモジニアス構成，ヘテロジニアス構成．}より多くの要求を処理しなければ
ならないサービスは，したがって多数のマシン上の多数のプロセスとして動作する．
これらのマシンがサービスの\emph{ノード}である．

\begin{definition}[スケールアウト]
  ノードをより強力なマシンで置き換える（\emph{スケールアップ}）のでは
  なく，ノードを追加することでサービスの処理能力を高める構成を
  \term[すけーるあうと]{スケールアウト}[scale-out]構成と呼ぶ．
\end{definition}

スケールアウト構成は，第3章のボス・ワーカーパターン%
\index{ぼすわーかーぱたーん@ボス・ワーカーパターン}をスレッドではなくマシン
に適用することが多い．フロントエンドのノードがすべての要求を受け取り，その
各々を，要求を処理するノードの1つへ渡す．

\begin{definition}[ロードバランサ]
  \term[ろーどばらんさ]{ロードバランサ}[load balancer]とは，サービスに
  到着する要求を，それを処理するノードの間へ分配するフロントエンドである．
\end{definition}

ボス・ワーカーパターンのワーカーがすべて等しいか特化しているかのいずれかで
あるのと同様に，スケールアウト構成のノードは2通りに編成できる
（\cref{fig:l16-scaleout}）．
\begin{description}
  \item[すべて等しい] どのノードも，あらゆる要求の処理のあらゆる段階を実行
    できる．ノードは機能的にホモジニアスである．
  \item[特化する] 各ノードは処理の一部の段階だけ，あるいは一部の種類の要求
    だけを実行する．ノードは機能的にヘテロジニアスである．
\end{description}

\begin{figure}[htbp]
  \centering
  % Scale-out architectures with a front-end load balancer: (a) functionally
% homogeneous nodes, (b) functionally heterogeneous (specialized) nodes.
% Usage: % Scale-out architectures with a front-end load balancer: (a) functionally
% homogeneous nodes, (b) functionally heterogeneous (specialized) nodes.
% Usage: % Scale-out architectures with a front-end load balancer: (a) functionally
% homogeneous nodes, (b) functionally heterogeneous (specialized) nodes.
% Usage: \input{figures/l16-scaleout}   (width ~118 mm)
\begin{tikzpicture}[x=1mm, y=1mm, font=\scriptsize\sffamily,
  >={Stealth[length=1.5mm]},
  fe/.style={draw, fill=ln-accent!18, minimum width=30mm, minimum height=8mm, inner sep=1pt, align=center},
  nd/.style={draw, fill=ln-box, minimum width=12mm, minimum height=9mm, inner sep=0.5pt, align=center},
  data/.style={draw, fill=ln-accent!35, minimum height=6mm, inner sep=1pt, align=center},
  lab/.style={font=\scriptsize\sffamily, text=ln-muted, align=center},
  ttl/.style={font=\footnotesize\sffamily\bfseries, anchor=west}]
  % ---------------- (a) homogeneous
  \begin{scope}[shift={(0,0)}]
    \node[ttl] at (0,56) {\iflangja{(a) ホモジニアス構成}{(a) homogeneous}};
    \draw[->, thick] (28.5,51) -- node[lab, right] {\iflangja{要求}{requests}} (28.5,44.5);
    \node[fe] (f) at (28.5,40) {\iflangja{フロントエンド}{front end}\\\textcolor{ln-muted}{\iflangja{要求を順に割り当てる}{assigns in turn}}};
    \foreach \i/\x in {1/7.5,2/21.5,3/35.5,4/49.5} {
      \node[nd] (n\i) at (\x,22) {\iflangja{ノード}{node}\\\textcolor{ln-muted}{A, B, C}};
      \draw[->] (f.south) -- (n\i.north);
      \draw (n\i.south) -- (\x,10);
    }
    \node[data, minimum width=54mm] at (28.5,7) {\iflangja{データ（全ノードからアクセス可能）}{data (accessible from every node)}};
  \end{scope}
  % ---------------- (b) heterogeneous
  \begin{scope}[shift={(61,0)}]
    \node[ttl] at (0,56) {\iflangja{(b) ヘテロジニアス構成}{(b) heterogeneous}};
    \draw[->, thick] (28.5,51) -- node[lab, right] {\iflangja{要求}{requests}} (28.5,44.5);
    \node[fe] (f) at (28.5,40) {\iflangja{フロントエンド}{front end}\\\textcolor{ln-muted}{\iflangja{要求の種類を調べる}{inspects each request}}};
    \foreach \i/\x/\t in {1/7.5/A,2/21.5/A,3/35.5/B,4/49.5/C} {
      \node[nd] (n\i) at (\x,22) {\iflangja{ノード}{node}\\\textcolor{ln-muted}{\t}};
      \draw[->] (f.south) -- (n\i.north);
      \draw (n\i.south) -- (\x,10);
    }
    \node[data, minimum width=26mm] at (14.5,7) {\iflangja{データ}{data} A};
    \node[data, minimum width=12mm] at (35.5,7) {B};
    \node[data, minimum width=12mm] at (49.5,7) {C};
  \end{scope}
  \node[lab, anchor=north] at (59,1.5)
    {\iflangja{A，B，C：要求の種類}{A, B, C: request types}};
\end{tikzpicture}
   (width ~118 mm)
\begin{tikzpicture}[x=1mm, y=1mm, font=\scriptsize\sffamily,
  >={Stealth[length=1.5mm]},
  fe/.style={draw, fill=ln-accent!18, minimum width=30mm, minimum height=8mm, inner sep=1pt, align=center},
  nd/.style={draw, fill=ln-box, minimum width=12mm, minimum height=9mm, inner sep=0.5pt, align=center},
  data/.style={draw, fill=ln-accent!35, minimum height=6mm, inner sep=1pt, align=center},
  lab/.style={font=\scriptsize\sffamily, text=ln-muted, align=center},
  ttl/.style={font=\footnotesize\sffamily\bfseries, anchor=west}]
  % ---------------- (a) homogeneous
  \begin{scope}[shift={(0,0)}]
    \node[ttl] at (0,56) {\iflangja{(a) ホモジニアス構成}{(a) homogeneous}};
    \draw[->, thick] (28.5,51) -- node[lab, right] {\iflangja{要求}{requests}} (28.5,44.5);
    \node[fe] (f) at (28.5,40) {\iflangja{フロントエンド}{front end}\\\textcolor{ln-muted}{\iflangja{要求を順に割り当てる}{assigns in turn}}};
    \foreach \i/\x in {1/7.5,2/21.5,3/35.5,4/49.5} {
      \node[nd] (n\i) at (\x,22) {\iflangja{ノード}{node}\\\textcolor{ln-muted}{A, B, C}};
      \draw[->] (f.south) -- (n\i.north);
      \draw (n\i.south) -- (\x,10);
    }
    \node[data, minimum width=54mm] at (28.5,7) {\iflangja{データ（全ノードからアクセス可能）}{data (accessible from every node)}};
  \end{scope}
  % ---------------- (b) heterogeneous
  \begin{scope}[shift={(61,0)}]
    \node[ttl] at (0,56) {\iflangja{(b) ヘテロジニアス構成}{(b) heterogeneous}};
    \draw[->, thick] (28.5,51) -- node[lab, right] {\iflangja{要求}{requests}} (28.5,44.5);
    \node[fe] (f) at (28.5,40) {\iflangja{フロントエンド}{front end}\\\textcolor{ln-muted}{\iflangja{要求の種類を調べる}{inspects each request}}};
    \foreach \i/\x/\t in {1/7.5/A,2/21.5/A,3/35.5/B,4/49.5/C} {
      \node[nd] (n\i) at (\x,22) {\iflangja{ノード}{node}\\\textcolor{ln-muted}{\t}};
      \draw[->] (f.south) -- (n\i.north);
      \draw (n\i.south) -- (\x,10);
    }
    \node[data, minimum width=26mm] at (14.5,7) {\iflangja{データ}{data} A};
    \node[data, minimum width=12mm] at (35.5,7) {B};
    \node[data, minimum width=12mm] at (49.5,7) {C};
  \end{scope}
  \node[lab, anchor=north] at (59,1.5)
    {\iflangja{A，B，C：要求の種類}{A, B, C: request types}};
\end{tikzpicture}
   (width ~118 mm)
\begin{tikzpicture}[x=1mm, y=1mm, font=\scriptsize\sffamily,
  >={Stealth[length=1.5mm]},
  fe/.style={draw, fill=ln-accent!18, minimum width=30mm, minimum height=8mm, inner sep=1pt, align=center},
  nd/.style={draw, fill=ln-box, minimum width=12mm, minimum height=9mm, inner sep=0.5pt, align=center},
  data/.style={draw, fill=ln-accent!35, minimum height=6mm, inner sep=1pt, align=center},
  lab/.style={font=\scriptsize\sffamily, text=ln-muted, align=center},
  ttl/.style={font=\footnotesize\sffamily\bfseries, anchor=west}]
  % ---------------- (a) homogeneous
  \begin{scope}[shift={(0,0)}]
    \node[ttl] at (0,56) {\iflangja{(a) ホモジニアス構成}{(a) homogeneous}};
    \draw[->, thick] (28.5,51) -- node[lab, right] {\iflangja{要求}{requests}} (28.5,44.5);
    \node[fe] (f) at (28.5,40) {\iflangja{フロントエンド}{front end}\\\textcolor{ln-muted}{\iflangja{要求を順に割り当てる}{assigns in turn}}};
    \foreach \i/\x in {1/7.5,2/21.5,3/35.5,4/49.5} {
      \node[nd] (n\i) at (\x,22) {\iflangja{ノード}{node}\\\textcolor{ln-muted}{A, B, C}};
      \draw[->] (f.south) -- (n\i.north);
      \draw (n\i.south) -- (\x,10);
    }
    \node[data, minimum width=54mm] at (28.5,7) {\iflangja{データ（全ノードからアクセス可能）}{data (accessible from every node)}};
  \end{scope}
  % ---------------- (b) heterogeneous
  \begin{scope}[shift={(61,0)}]
    \node[ttl] at (0,56) {\iflangja{(b) ヘテロジニアス構成}{(b) heterogeneous}};
    \draw[->, thick] (28.5,51) -- node[lab, right] {\iflangja{要求}{requests}} (28.5,44.5);
    \node[fe] (f) at (28.5,40) {\iflangja{フロントエンド}{front end}\\\textcolor{ln-muted}{\iflangja{要求の種類を調べる}{inspects each request}}};
    \foreach \i/\x/\t in {1/7.5/A,2/21.5/A,3/35.5/B,4/49.5/C} {
      \node[nd] (n\i) at (\x,22) {\iflangja{ノード}{node}\\\textcolor{ln-muted}{\t}};
      \draw[->] (f.south) -- (n\i.north);
      \draw (n\i.south) -- (\x,10);
    }
    \node[data, minimum width=26mm] at (14.5,7) {\iflangja{データ}{data} A};
    \node[data, minimum width=12mm] at (35.5,7) {B};
    \node[data, minimum width=12mm] at (49.5,7) {C};
  \end{scope}
  \node[lab, anchor=north] at (59,1.5)
    {\iflangja{A，B，C：要求の種類}{A, B, C: request types}};
\end{tikzpicture}

  \caption{スケールアウト構成．(a) どのノードもあらゆる種類の要求を処理し，
    どのノードもすべてのデータにアクセスできる．(b) ノードは要求の種類に
    特化しており，フロントエンドはノードを選ぶために各要求を調べなければ
    ならない．}
  \label{fig:l16-scaleout}
\end{figure}

\subsection{ホモジニアス構成}

\begin{definition}[ホモジニアス構成]
  \term[ほもじにあすこうせい]{ホモジニアス構成}[homogeneous architecture]
  では，すべてのノードが機能的に同一である．どのノードも，あらゆる要求に
  対するあらゆる処理段階を実行できる．
\end{definition}

ホモジニアス構成のフロントエンドは単純である．要求を調べる必要はなく，
ラウンドロビンで順に，あるいは最も負荷の小さいノードへ割り当てればよい．
機能的にホモジニアスであっても，すべてのノードがすべてのデータを格納して
いる必要はない．必要なのは，どのノードもすべてのデータに\emph{アクセス
できる}ことだけであり，たとえば分散ファイルシステムや共有のデータベース層を
通じてであってよい．その層自体は複製されていても分割されていてもよい
（第14章）．処理能力を高めるには，単にノードを追加する．

その代価は，サービスがキャッシュからほとんど利益を得られないことである．
フロントエンドは，どのノードが最近似た要求を処理し，必要なデータをホット
キャッシュ\index{ほっときゃっしゅ@ホットキャッシュ}に保持しているかを
知らない．あるノードがキャッシュから処理できたはずの要求も，したがって，
まずデータを取得しなければならないノードへ送られる可能性が高い．これは，
等しいワーカーについて第3章が指摘した局所性の欠如である．

\subsection{ヘテロジニアス構成}

\begin{definition}[ヘテロジニアス構成]
  \term[へてろじにあすこうせい]{ヘテロジニアス構成}%
  [heterogeneous architecture]では，ノードが機能的に特化している．異なる
  ノードが異なる処理段階を実行し，あるいは異なる種類の要求を処理する．
\end{definition}

このときデータは，すべてのノードから一様にアクセスできる必要はない．各
ノードには，自分の処理段階や要求の種類に必要なデータだけがあればよい．
利点は局所性である．ある種類の要求はすべて同じ少数のノードに到達し，それら
のノードは，第3章の特化したスレッド\index{とっか@特化}と同様に，その要求が
必要とするコードとデータをホットキャッシュに保つ．しかしフロントエンドは
より複雑になる．\margin{接続を転送するだけのフロントエンドは，トランスポート層
のヘッダより先を見なくてよい．URL や cookie でノードを選ぶものは，接続を終端
して要求そのものを解析しなければならない．}どのノードがその要求を処理できるか
を判定するために，たとえば URL を解析して，あらゆる要求を調べなければならない．
サービスの管理もまた複雑になる．

\subsection{スケーリングと管理}

2つの構成の違いは，負荷が増えたときに最も大きく現れる．\margin{スケール
アウトとスケールアップは，水平スケーリングおよび垂直スケーリングとも呼ば
れる．いずれもサービスのスループットを高めるが，どちらも単一の要求の処理を
短くはしない．}ホモジニアス構成をスケールさせるには，ノードと，それが必要と
するプロセス，マシン，ストレージへのアクセスを追加すればよい．新しいノードは
どれもあらゆる種類の要求を分担するので，どの要求が負荷を生んでいるかを知る
必要はない．ヘテロジニアス構成をスケールさせるには，運用者は
どの種類の要求に需要があるかを知り，その種類のノードを追加しなければならず，
そのためにはワークロード\index{わーくろーど@ワークロード}のプロファイリング
が必要である．要求の種類ごとの需要は時とともに変化するので，ノードの構成比も
それに追随しなければならない．さもなければ，人気のある種類のノードは未処理の
要求の長い待ち行列を抱えた\emph{ホットスポット}となり，他の種類のノードは
遊休する．\Cref{tab:l16-arch} にこの比較をまとめる．

\begin{table}[htbp]
  \centering
  \caption{ホモジニアスなスケールアウト構成とヘテロジニアスなスケールアウト
    構成．}
  \label{tab:l16-arch}
  \small
  \begin{tabular}{@{}>{\raggedright\arraybackslash}p{0.2\linewidth}>{\raggedright\arraybackslash}p{0.36\linewidth}>{\raggedright\arraybackslash}p{0.36\linewidth}@{}}
    \toprule
    & ホモジニアス & ヘテロジニアス \\
    \midrule
    ノード & 機能的に同一 & 処理段階または要求の種類に特化 \\
    フロントエンド & 単純；要求を調べる必要がない & あらゆる要求を調べ
      なければならない \\
    データ & どのノードからもアクセスできる & 一部のノードにだけ必要 \\
    キャッシュと局所性 & ほとんど利益がない & ある種類の要求が同じノードに
      到達する \\
    スケーリング & ノードを追加する & 需要をプロファイルし，必要な種類の
      ノードを追加する \\
    要求の構成比の変化 & 介入なしに吸収される & ノードを割り当て直さない
      限りホットスポットが生じる \\
    \bottomrule
  \end{tabular}
\end{table}

\begin{example}
  \label{ex:l16-scaling}
  あるオンラインショップが2種類の要求を受け取る．ホモジニアス構成のノード
  上では，検索1件の処理に \SI{5}{\milli\second}，チェックアウト1件に
  \SI{20}{\milli\second} かかる．毎秒1200件の検索と200件のチェックアウトと
  いう負荷は，1秒当たり $1200 \times 0.005 + 200 \times 0.020 = 6 + 4 = 10$
  秒の処理を要し，これは10ノード分の処理能力である．構成比が毎秒600件の検索
  と350件のチェックアウトに変化しても，必要量は $3 + 7 = 10$ で変わらない．
  同じ10ノードで足り，この変化に誰も気づく必要がない．

  ヘテロジニアス構成では，ホットキャッシュによって，特化したノードが要求を
  \SI{20}{\percent} 短い時間で，すなわち検索1件を \SI{4}{\milli\second}，
  チェックアウト1件を \SI{16}{\milli\second} で処理できるとする．元の負荷に
  は $\lceil 1200 \times 0.004 \rceil = 5$ 個の検索ノードと
  $\lceil 200 \times 0.016 \rceil = 4$ 個のチェックアウトノード，合計9個が
  必要である．変化の後には $\lceil 2.4 \rceil = 3$ 個の検索ノードと
  $\lceil 5.6 \rceil = 6$ 個のチェックアウトノードが必要となる．合計は同じ
  9個だが，構成比は異なる．運用者が2個のノードを検索からチェックアウトへ
  移すまで，4個のチェックアウトノードは毎秒 $4/0.016 = 250$ 件しか完了でき
  ず，未処理の待ち行列は毎秒100件ずつ伸び，5個の検索ノードのうち2個は遊休
  する．
\end{example}

\begin{keypoint}
  インターネットサービスは層に分けて構成され，ロードバランサの背後で多数の
  ノードへスケールアウトされる．ホモジニアスなノードはフロントエンドと
  スケーリングを単純に保つが，キャッシュからほとんど利益を得られない．
  ヘテロジニアスなノードは局所性を得るが，その代価として，要求を調べる
  フロントエンドと，要求の構成比に追随する管理を必要とする．
\end{keypoint}

\section{クラウドコンピューティングの動機}
\label{sec:l16-motivation}

オンライン小売業者である Amazon は，ホリデーシーズンの商戦のピーク負荷に
合わせてハードウェアを用意しており，1年の残りの期間，そのハードウェアの
多くは遊休していた．\source{P4L4，クラウドコンピューティングの代表例：
Animoto，クラウドコンピューティングの要件；Armbrust ら
\cite{armbrust2009}．}Amazon は Web ベースの API を通じてそれへのアクセスを
開放し，第三者が料金を払って Amazon のハードウェア上で自分のワークロードを
実行できるようにした．これが Amazon Web Services（AWS）である．2006年に開始
されたその Elastic Compute Cloud（EC2）は\emph{コンピュートインスタンス}，
すなわち顧客が API を通じて起動・停止し，時間単位で料金を支払う仮想マシン
（第12章）を貸し出す．\margin{EC2 は2006年8月に限定的な
パブリックベータとして公開され，2008年10月に一般提供となった．}

利用者の写真から動画を生成するサービスである Animoto は，コンピュート
インスタンスを EC2 から借りていた．2008年4月に Animoto が Facebook の利用者に
も使えるようになると需要は急増し，3日のうちに利用するサーバ数は50台から3500台
へ増え，その後トラフィックはピークをはるかに下回る水準まで下がった
\cite{armbrust2009}．数日で2桁近く成長することは，マシンを購入し，設置し，
構成しなければならない場合には不可能である．借りた資源では，必要なのは API の
呼び出しだけであった．

\subsection{変動する需要に対する資源の用意}

従来は，サービスは自分の資源を購入し構成する．その容量は期待される需要，
通常はその期待されるピークから事前に決められ，後から容量を増やせるのは大きな
刻みに限られ，そのたびに購入による遅れを伴う．需要は時とともに変動するので，
固定された容量はほとんどの時点で適切でない（\cref{fig:l16-capacity}a）．
\begin{itemize}
  \item 需要が容量を下回っている間，遊休する資源は無駄になる．支払い済みで
    ありながら使われていない．
  \item 需要が容量を超えると，要求は捨てられ，それらの利用者に提供する機会
    と，それがもたらしたはずの収益が失われる．
\end{itemize}

\begin{figure}[htbp]
  \centering
  % Capacity versus demand over time: (a) fixed capacity provisioned for the
% expected peak wastes resources and loses requests during an unexpected
% surge; (b) elastic capacity follows the demand.
% Usage: % Capacity versus demand over time: (a) fixed capacity provisioned for the
% expected peak wastes resources and loses requests during an unexpected
% surge; (b) elastic capacity follows the demand.
% Usage: % Capacity versus demand over time: (a) fixed capacity provisioned for the
% expected peak wastes resources and loses requests during an unexpected
% surge; (b) elastic capacity follows the demand.
% Usage: \input{figures/l16-capacity}   (width ~118 mm)
\begin{tikzpicture}[x=1mm, y=1mm, font=\scriptsize\sffamily,
  >={Stealth[length=1.5mm]},
  ax/.style={->, thin},
  dem/.style={thick, ln-accent},
  capa/.style={thick, ln-caution},
  lab/.style={font=\scriptsize\sffamily, align=center},
  ttl/.style={font=\footnotesize\sffamily\bfseries, anchor=west}]
  % ---------------- (a) fixed capacity
  \begin{scope}[shift={(3,0)}]
    \node[ttl] at (-4,43) {\iflangja{(a) 固定された容量}{(a) fixed capacity}};
    \begin{scope}
      \clip (0,0) rectangle (50,22);
      \fill[ln-rule!45] (0,34) -- plot[smooth, tension=0.6] coordinates
        {(0,6) (5,8) (10,12) (15,10) (20,14) (25,17) (30,15) (35,19) (38,29) (41,28) (44,18) (50,16)}
        -- (50,34) -- cycle;
    \end{scope}
    \begin{scope}
      \clip (0,22) rectangle (50,34);
      \fill[ln-caution!30] (0,0) -- plot[smooth, tension=0.6] coordinates
        {(0,6) (5,8) (10,12) (15,10) (20,14) (25,17) (30,15) (35,19) (38,29) (41,28) (44,18) (50,16)}
        -- (50,0) -- cycle;
    \end{scope}
    \draw[dem] plot[smooth, tension=0.6] coordinates
      {(0,6) (5,8) (10,12) (15,10) (20,14) (25,17) (30,15) (35,19) (38,29) (41,28) (44,18) (50,16)};
    \draw[capa] (0,22) -- (50,22);
    \draw[ax] (0,0) -- (53,0) node[below left, lab] {\iflangja{時間}{time}};
    \draw[ax] (0,0) -- (0,35) node[above right, lab, anchor=south west] {\iflangja{資源}{resources}};
    \node[lab, fill=white, inner sep=0.6pt] at (14,18.5) {\iflangja{無駄な資源}{wasted resources}};
    \draw[->, ln-caution] (30,30) -- (37,25.5);
    \node[lab, anchor=east, text=ln-caution] at (30.5,30.5) {\iflangja{逸失した機会}{lost opportunity}};
  \end{scope}
  % ---------------- (b) elastic capacity
  \begin{scope}[shift={(63,0)}]
    \node[ttl] at (-4,43) {\iflangja{(b) 伸縮する容量}{(b) elastic capacity}};
    \draw[dem] plot[smooth, tension=0.6] coordinates
      {(0,6) (5,8) (10,12) (15,10) (20,14) (25,17) (30,15) (35,19) (38,29) (41,28) (44,18) (50,16)};
    \draw[capa] (0,9) -- (5,9) -- (5,13.5) -- (10,13.5) -- (10,14) -- (15,14)
      -- (15,15.5) -- (20,15.5) -- (20,19) -- (30,19) -- (30,21) -- (35,21)
      -- (35,31) -- (43,31) -- (43,21) -- (46,21) -- (46,18.5) -- (50,18.5);
    \draw[ax] (0,0) -- (53,0) node[below left, lab] {\iflangja{時間}{time}};
    \draw[ax] (0,0) -- (0,35) node[above right, lab, anchor=south west] {\iflangja{資源}{resources}};
  \end{scope}
  % ---------------- legend
  \draw[dem] (36,-7) -- (44,-7);
  \node[lab, anchor=west] at (45,-7) {\iflangja{需要}{demand}};
  \draw[capa] (62,-7) -- (70,-7);
  \node[lab, anchor=west] at (71,-7) {\iflangja{容量}{capacity}};
\end{tikzpicture}
   (width ~118 mm)
\begin{tikzpicture}[x=1mm, y=1mm, font=\scriptsize\sffamily,
  >={Stealth[length=1.5mm]},
  ax/.style={->, thin},
  dem/.style={thick, ln-accent},
  capa/.style={thick, ln-caution},
  lab/.style={font=\scriptsize\sffamily, align=center},
  ttl/.style={font=\footnotesize\sffamily\bfseries, anchor=west}]
  % ---------------- (a) fixed capacity
  \begin{scope}[shift={(3,0)}]
    \node[ttl] at (-4,43) {\iflangja{(a) 固定された容量}{(a) fixed capacity}};
    \begin{scope}
      \clip (0,0) rectangle (50,22);
      \fill[ln-rule!45] (0,34) -- plot[smooth, tension=0.6] coordinates
        {(0,6) (5,8) (10,12) (15,10) (20,14) (25,17) (30,15) (35,19) (38,29) (41,28) (44,18) (50,16)}
        -- (50,34) -- cycle;
    \end{scope}
    \begin{scope}
      \clip (0,22) rectangle (50,34);
      \fill[ln-caution!30] (0,0) -- plot[smooth, tension=0.6] coordinates
        {(0,6) (5,8) (10,12) (15,10) (20,14) (25,17) (30,15) (35,19) (38,29) (41,28) (44,18) (50,16)}
        -- (50,0) -- cycle;
    \end{scope}
    \draw[dem] plot[smooth, tension=0.6] coordinates
      {(0,6) (5,8) (10,12) (15,10) (20,14) (25,17) (30,15) (35,19) (38,29) (41,28) (44,18) (50,16)};
    \draw[capa] (0,22) -- (50,22);
    \draw[ax] (0,0) -- (53,0) node[below left, lab] {\iflangja{時間}{time}};
    \draw[ax] (0,0) -- (0,35) node[above right, lab, anchor=south west] {\iflangja{資源}{resources}};
    \node[lab, fill=white, inner sep=0.6pt] at (14,18.5) {\iflangja{無駄な資源}{wasted resources}};
    \draw[->, ln-caution] (30,30) -- (37,25.5);
    \node[lab, anchor=east, text=ln-caution] at (30.5,30.5) {\iflangja{逸失した機会}{lost opportunity}};
  \end{scope}
  % ---------------- (b) elastic capacity
  \begin{scope}[shift={(63,0)}]
    \node[ttl] at (-4,43) {\iflangja{(b) 伸縮する容量}{(b) elastic capacity}};
    \draw[dem] plot[smooth, tension=0.6] coordinates
      {(0,6) (5,8) (10,12) (15,10) (20,14) (25,17) (30,15) (35,19) (38,29) (41,28) (44,18) (50,16)};
    \draw[capa] (0,9) -- (5,9) -- (5,13.5) -- (10,13.5) -- (10,14) -- (15,14)
      -- (15,15.5) -- (20,15.5) -- (20,19) -- (30,19) -- (30,21) -- (35,21)
      -- (35,31) -- (43,31) -- (43,21) -- (46,21) -- (46,18.5) -- (50,18.5);
    \draw[ax] (0,0) -- (53,0) node[below left, lab] {\iflangja{時間}{time}};
    \draw[ax] (0,0) -- (0,35) node[above right, lab, anchor=south west] {\iflangja{資源}{resources}};
  \end{scope}
  % ---------------- legend
  \draw[dem] (36,-7) -- (44,-7);
  \node[lab, anchor=west] at (45,-7) {\iflangja{需要}{demand}};
  \draw[capa] (62,-7) -- (70,-7);
  \node[lab, anchor=west] at (71,-7) {\iflangja{容量}{capacity}};
\end{tikzpicture}
   (width ~118 mm)
\begin{tikzpicture}[x=1mm, y=1mm, font=\scriptsize\sffamily,
  >={Stealth[length=1.5mm]},
  ax/.style={->, thin},
  dem/.style={thick, ln-accent},
  capa/.style={thick, ln-caution},
  lab/.style={font=\scriptsize\sffamily, align=center},
  ttl/.style={font=\footnotesize\sffamily\bfseries, anchor=west}]
  % ---------------- (a) fixed capacity
  \begin{scope}[shift={(3,0)}]
    \node[ttl] at (-4,43) {\iflangja{(a) 固定された容量}{(a) fixed capacity}};
    \begin{scope}
      \clip (0,0) rectangle (50,22);
      \fill[ln-rule!45] (0,34) -- plot[smooth, tension=0.6] coordinates
        {(0,6) (5,8) (10,12) (15,10) (20,14) (25,17) (30,15) (35,19) (38,29) (41,28) (44,18) (50,16)}
        -- (50,34) -- cycle;
    \end{scope}
    \begin{scope}
      \clip (0,22) rectangle (50,34);
      \fill[ln-caution!30] (0,0) -- plot[smooth, tension=0.6] coordinates
        {(0,6) (5,8) (10,12) (15,10) (20,14) (25,17) (30,15) (35,19) (38,29) (41,28) (44,18) (50,16)}
        -- (50,0) -- cycle;
    \end{scope}
    \draw[dem] plot[smooth, tension=0.6] coordinates
      {(0,6) (5,8) (10,12) (15,10) (20,14) (25,17) (30,15) (35,19) (38,29) (41,28) (44,18) (50,16)};
    \draw[capa] (0,22) -- (50,22);
    \draw[ax] (0,0) -- (53,0) node[below left, lab] {\iflangja{時間}{time}};
    \draw[ax] (0,0) -- (0,35) node[above right, lab, anchor=south west] {\iflangja{資源}{resources}};
    \node[lab, fill=white, inner sep=0.6pt] at (14,18.5) {\iflangja{無駄な資源}{wasted resources}};
    \draw[->, ln-caution] (30,30) -- (37,25.5);
    \node[lab, anchor=east, text=ln-caution] at (30.5,30.5) {\iflangja{逸失した機会}{lost opportunity}};
  \end{scope}
  % ---------------- (b) elastic capacity
  \begin{scope}[shift={(63,0)}]
    \node[ttl] at (-4,43) {\iflangja{(b) 伸縮する容量}{(b) elastic capacity}};
    \draw[dem] plot[smooth, tension=0.6] coordinates
      {(0,6) (5,8) (10,12) (15,10) (20,14) (25,17) (30,15) (35,19) (38,29) (41,28) (44,18) (50,16)};
    \draw[capa] (0,9) -- (5,9) -- (5,13.5) -- (10,13.5) -- (10,14) -- (15,14)
      -- (15,15.5) -- (20,15.5) -- (20,19) -- (30,19) -- (30,21) -- (35,21)
      -- (35,31) -- (43,31) -- (43,21) -- (46,21) -- (46,18.5) -- (50,18.5);
    \draw[ax] (0,0) -- (53,0) node[below left, lab] {\iflangja{時間}{time}};
    \draw[ax] (0,0) -- (0,35) node[above right, lab, anchor=south west] {\iflangja{資源}{resources}};
  \end{scope}
  % ---------------- legend
  \draw[dem] (36,-7) -- (44,-7);
  \node[lab, anchor=west] at (45,-7) {\iflangja{需要}{demand}};
  \draw[capa] (62,-7) -- (70,-7);
  \node[lab, anchor=west] at (71,-7) {\iflangja{容量}{capacity}};
\end{tikzpicture}

  \caption{時間に対する容量と需要．(a) 固定された容量は，需要が低い間は資源
    を無駄にし，需要が急増する間は要求を失う．(b) 伸縮する容量は需要に追随
    する．}
  \label{fig:l16-capacity}
\end{figure}

理想的には，サービスの容量は需要に追随するのが望ましい
（\cref{fig:l16-capacity}b）．
\begin{itemize}
  \item 容量は需要に応じて，増える方向にも減る方向にも伸縮する．
  \item スケーリングは即座に行われる．
  \item 費用は需要に，したがって収益の機会に比例する．
  \item 以上のすべてが，特別な専門知識なしに自動的に行われる．
  \item 資源にはいつでも，どこからでもアクセスできる．
\end{itemize}

\begin{definition}[伸縮性]
  需要の増減に応じてシステムが資源を素早く自動的に獲得し解放し，その容量が
  需要に追随する能力を\term[しんしゅくせい]{伸縮性}[elasticity]と呼ぶ．
\end{definition}

この理想の代価は，サービスが，自分の動作する基盤である資源をもはや所有せず，
そのデータが他者の運用するマシン上に置かれることである．\caution{容量を借りて
も，資源をどれだけ用意するかという判断がなくなるわけではない．その判断の時間軸
が短くなるだけである．いくつのインスタンスを動かすかは依然として決めなければ
ならず，それが数年ごとではなく数分ごとになる．}

\begin{example}
  あるサービスの需要が，毎日16時間はサーバ20台，残りの8時間はサーバ100台
  であるとする．1日当たり $20 \times 16 + 100 \times 8 = 1120$ サーバ時で
  ある．所有するサーバは，使っていても使っていなくても1時間当たり $c$ の
  費用がかかる．ピークに合わせて用意すると，サービスは
  $100 \times 24 = 2400$ サーバ時，すなわち1日当たり $2400c$ を支払い，その
  うち $1 - 1120/2400 \approx \SI{53}{\percent}$ が無駄になる．平均の需要で
  ある $1120/24 \approx 47$ 台に合わせて用意しても，需要が低い間に
  $27 \times 16 = 432$ サーバ時を無駄にし，ピークの間は要求の
  $53/100 = \SI{53}{\percent}$ を捨てる．\SI{50}{\percent} 高い価格，すなわち
  1サーバ時当たり $1.5c$ で伸縮する容量を借りると，サービスの支払いは1日
  当たり $1120 \times 1.5c = 1680c$ となり，ピークに合わせて用意する場合より
  \SI{30}{\percent} 少なく，要求を1つも捨てない．
\end{example}

\section{クラウドコンピューティング}
\label{sec:l16-cloud}

クラウドコンピューティングは，上で述べた理想に近づくものである．
\source{P4L4，クラウドコンピューティングの概観，クラウドコンピューティングが
成り立つ理由，クラウドコンピューティングの構想，クラウドコンピューティングの
定義；Armbrust ら \cite{armbrust2009}；Mell と Grance \cite{mell2011}；
Silberschatz ら ch.~1．}その本質的な要素は次のとおりである．
\begin{description}
  \item[共有される資源] プロバイダは，計算，ストレージ，ネットワークといった
    インフラストラクチャに加えて，ソフトウェアやサービスを提供し，それらを
    多数の顧客の間で共有する．
  \item[アクセスと設定のための API] 顧客は，Web ベースの API，プログラミング
    言語のライブラリ，あるいはコマンドラインツールを通じて資源を要求し，
    設定し，解放する．
  \item[課金と会計] 利用量は計測され課金される．多様な価格モデルが併存する．
    たとえば，オンデマンド価格，長期間予約した容量に対するより低い価格，
    そして需給に応じて変動する余剰容量の\emph{スポット}価格であり，プロ
    バイダは事実上の市場を運営することになる．資源は離散的な大きさで提供
    される．たとえば CPU 数とメモリ量の異なる small，medium，extra-large と
    いったインスタンスである．
  \item[プロバイダによる管理] プロバイダが，データセンタ，ハードウェア，
    およびサービスを実装するソフトウェアを運用する．
\end{description}

\begin{definition}[クラウドコンピューティング]
  多数の顧客の間で共有される資源から，API を通じて需要に応じて獲得・解放
  され，利用に応じて課金され，プロバイダによって管理される計算資源を，
  ネットワーク越しにサービスとして提供することを，
  \term[くらうどこんぴゅーてぃんぐ]{クラウドコンピューティング}%
  [cloud computing]と呼ぶ．
\end{definition}

\begin{note}
  公表されている定義もこれらの要素と一致する．NIST の定義 \cite{mell2011} は
  5つの本質的な特徴，すなわちオンデマンドのセルフサービス，幅広い
  ネットワークアクセス，資源のプール化，迅速な伸縮性，計測されたサービスを
  挙げる．Armbrust ら \cite{armbrust2009} はこの語を，インターネット越しに
  サービスとして提供されるアプリケーションと，それを提供するデータセンタ内の
  ハードウェアおよびシステムソフトウェアの両方に用いる．後者を彼らは
  \emph{クラウド}と呼ぶ．
\end{note}

\subsection{クラウドコンピューティングが成り立つ理由}

プロバイダは，すべての顧客の需要を賄うだけの資源を所有し，なお顧客が自前の
資源に費やすであろう額より安く課金しなければならない．これを可能にする原理が
2つある．

\begin{definition}[大数の法則]
  \term[たいすうのほうそく]{大数の法則}[law of large numbers]によれば，同じ
  分布に従う $N$ 個の独立な確率量の平均は，$N$ が大きくなるにつれてその
  期待値に近づく．
\end{definition}

単一の顧客の資源需要は時とともに大きく変動するが，異なる顧客はそれぞれ別の
時点でピークを迎えるので，多数の顧客にわたって平均した需要はほぼ一定である．
$N$ 人の顧客の需要が独立で，各々の平均が $\mu$，標準偏差が $\sigma$ であれ
ば，その和の平均は $N\mu$，標準偏差は $\sigma\sqrt{N}$ となり，平均に対する
変動の割合 $\sigma/(\mu\sqrt{N})$ は $N$ とともに小さくなる．\tip{相対的な
変動は $1/\sqrt{N}$ でしか減らない．独立な顧客が10倍になっても約3分の1に，
100倍になってようやく10分の1になるにすぎない．}したがってプロバイダは，個々
のピークの和よりはるかに小さい，合計の需要に合わせて資源を用意できる．

\begin{example}
  あるプロバイダが100の顧客にサービスを提供している．各顧客は普段10台の
  サーバを必要とし，ピーク時には100台を必要とする．ピークは時間全体の
  \SI{10}{\percent} を占め，他の顧客のピークとは独立に生じる．すべての顧客の
  ピークに合わせて用意すると $100 \times 100 = \num{10000}$ 台のサーバが要る．
  しかし，ある時点でピークにある顧客の数 $K$ は，平均が
  $100 \times 0.1 = 10$，標準偏差が $\sqrt{100 \times 0.1 \times 0.9} = 3$
  であり，合計の需要である $1000 + 90K$ 台は平均1900台となる．平均より標準
  偏差3つ分上の $K \le 19$ を賄える \num{2710} 台のサーバがあれば，需要が
  容量を超える確率はわずか \SI{0.2}{\percent} 程度である．プロバイダに必要な
  サーバは，顧客が自前で所有する場合の台数の \SI{30}{\percent} に満たない．
\end{example}

第2の原理は\term[きぼのけいざい]{規模の経済}[economies of scale]である．
資源を提供する単位当たりの費用は，提供する量が増えるにつれて下がる．
プロバイダは，ハードウェア，電力，ネットワーク容量をまとめて購入し，データ
センタの固定費を多数のマシンに分散し，管理者1人に多数のマシンを管理させる．

\subsection{ユーティリティとしてのコンピューティング}

クラウドコンピューティングの構想は，その技術より古い．1961年，John McCarthy
は，計算がいつか電話網のように公共のユーティリティとして組織されうること，
そしてそのような計算ユーティリティが重要な新産業の基盤となりうることを示唆
した．

\begin{definition}[ユーティリティコンピューティング]
  \term[ゆーてぃりてぃこんぴゅーてぃんぐ]{ユーティリティコンピューティング}%
  [utility computing]では，計算資源が電力や電話サービスと同じように提供
  される．すなわち，計測され，代替可能な---どの単位も他の単位と変わらない---%
  商品として提供され，顧客は必要に応じてそれを消費し，利用に応じて支払う．
\end{definition}

Armbrust ら \cite{armbrust2009} の用語では，ユーティリティコンピューティング
とは，クラウドが従量課金で一般に公開されるときに販売するものである．現在の
クラウドがこの構想を実現しているのは一部にとどまる．計算が完全に代替可能な
ユーティリティにはなっていないからである．その限界は次のとおりである．
\begin{description}
  \item[API のロックイン] プロバイダはそれぞれ独自の API を持ち，ある
    プロバイダ向けに書かれたアプリケーションは，変更なしに別のプロバイダへ
    移せない．
  \item[ハードウェアへの依存] アプリケーションは，特定のプロセッサや
    アクセラレータといった特定のハードウェアに依存することがあり，それを
    すべてのプロバイダやインスタンス種別が提供するわけではない．
  \item[レイテンシ] 資源は遠隔のデータセンタに存在し，それとのやり取りは
    すべてネットワークのレイテンシを伴う．
  \item[プライバシとセキュリティ] データと計算は，顧客が所有も制御もせず，
    他の顧客と共有されるマシン上に存在する．
\end{description}

\begin{keypoint}
  クラウドコンピューティングは，共有され，プロバイダによって管理される資源
  を，API を通じて伸縮的に，利用に応じた課金で提供する．これが成り立つのは，
  多数の顧客の合計の需要が個々の顧客の需要よりはるかに変動しないから（大数の
  法則）であり，また単位当たりの費用が規模とともに下がるから（規模の経済）
  である．ユーティリティとしての計算という構想を実現しているのは一部に
  とどまる．
\end{keypoint}

\section{デプロイメントモデルとサービスモデル}
\label{sec:l16-models}

クラウドは，誰が使えるか，そして何を提供するかによって異なる．\source{P4L4，
クラウドのデプロイメントモデル，クラウドのサービスモデル；Mell と Grance
\cite{mell2011}；Silberschatz ら ch.~1．}\emph{デプロイメントモデル}は，誰が
クラウドを使うかを表す．
\begin{description}
  \item[パブリッククラウド] \term[ぱぶりっくくらうど]{パブリッククラウド}%
    [public cloud]は，プロバイダが第三者の顧客，すなわちその\emph{テナント}
    のために運用し，テナントが資源を共有するクラウドである．
  \item[プライベートクラウド] \term[ぷらいべーとくらうど]{プライベート
    クラウド}[private cloud]では，組織が，仮想化，API，自動化された資源
    割り当てといったクラウドの技術を，自分のインフラストラクチャに，自分の
    アプリケーションのために適用する．
  \item[ハイブリッドクラウド] \term[はいぶりっどくらうど]{ハイブリッド
    クラウド}[hybrid cloud]は，プライベートクラウドとパブリッククラウドを
    組み合わせる．組織は自分のプライベートクラウドでアプリケーションを動かし，
    フェイルオーバのため，自分の容量を超える需要の急増を吸収するため，
    あるいはテストのためにパブリッククラウドを使う．
  \item[コミュニティクラウド] \term[こみゅにてぃくらうど]{コミュニティ
    クラウド}[community cloud]は，政府機関や研究機関など，ある種類の利用者に
    利用が限られたパブリッククラウドである．
\end{description}

\emph{サービスモデル}は，ハードウェアとソフトウェアのスタックのうちどこまで
をプロバイダが管理するかを表す（\cref{fig:l16-service-models}）．オンプレミス
の設置では，組織がスタック全体を自分で管理する．ネットワーク，ストレージ，
サーバから，仮想化，オペレーティングシステム，ミドルウェア，ランタイムを
経て，自分のデータとアプリケーションまでである．

\begin{figure}[htbp]
  \centering
  % Cloud service models: which layers of the software and hardware stack the
% customer manages and which the provider manages, for an on-premises
% installation, IaaS, PaaS and SaaS.
% Usage: % Cloud service models: which layers of the software and hardware stack the
% customer manages and which the provider manages, for an on-premises
% installation, IaaS, PaaS and SaaS.
% Usage: % Cloud service models: which layers of the software and hardware stack the
% customer manages and which the provider manages, for an on-premises
% installation, IaaS, PaaS and SaaS.
% Usage: \input{figures/l16-service-models}   (width ~116 mm)
\begin{tikzpicture}[x=1mm, y=1mm, font=\scriptsize\sffamily,
  own/.style={draw, fill=white, minimum width=26mm, minimum height=4.6mm, inner sep=0.5pt},
  prov/.style={draw, fill=ln-accent!25, minimum width=26mm, minimum height=4.6mm, inner sep=0.5pt},
  hd/.style={font=\footnotesize\sffamily\bfseries, anchor=base},
  ex/.style={font=\scriptsize\sffamily, text=ln-muted, anchor=south},
  lab/.style={font=\scriptsize\sffamily, anchor=west}]
  % layers from top (1) to bottom (9)
  \def\layername#1{%
    \ifcase#1\or \iflangja{アプリケーション}{applications}%
    \or \iflangja{データ}{data}%
    \or \iflangja{ランタイム}{runtime}%
    \or \iflangja{ミドルウェア}{middleware}%
    \or OS%
    \or \iflangja{仮想化}{virtualization}%
    \or \iflangja{サーバ}{servers}%
    \or \iflangja{ストレージ}{storage}%
    \or \iflangja{ネットワーク}{networking}\fi}
  % column / header / example / first layer managed by the provider (10 = none)
  \foreach \c/\h/\e/\p in {%
      0/{\iflangja{オンプレミス}{on premises}}/{}/10,
      1/IaaS/{Amazon EC2}/6,
      2/PaaS/{Google App Engine}/3,
      3/SaaS/{Gmail}/1} {
    \pgfmathsetmacro{\cx}{14 + 30*\c}
    \node[hd] at (\cx,52.3) {\h};
    \node[ex] at (\cx,46.8) {\strut\e};
    \foreach \l in {1,...,9} {
      \pgfmathsetmacro{\ly}{46 - 4.6*\l + 2.3}
      \ifnum\l<\p
        \node[own] at (\cx,\ly) {\layername{\l}};
      \else
        \node[prov] at (\cx,\ly) {\layername{\l}};
      \fi
    }
  }
  % legend
  \node[own, minimum width=5mm, minimum height=3.5mm] at (27,-3.5) {};
  \node[lab] at (30,-3.5) {\iflangja{利用者が管理}{managed by the customer}};
  \node[prov, minimum width=5mm, minimum height=3.5mm] at (70,-3.5) {};
  \node[lab] at (73,-3.5) {\iflangja{プロバイダが管理}{managed by the provider}};
\end{tikzpicture}
   (width ~116 mm)
\begin{tikzpicture}[x=1mm, y=1mm, font=\scriptsize\sffamily,
  own/.style={draw, fill=white, minimum width=26mm, minimum height=4.6mm, inner sep=0.5pt},
  prov/.style={draw, fill=ln-accent!25, minimum width=26mm, minimum height=4.6mm, inner sep=0.5pt},
  hd/.style={font=\footnotesize\sffamily\bfseries, anchor=base},
  ex/.style={font=\scriptsize\sffamily, text=ln-muted, anchor=south},
  lab/.style={font=\scriptsize\sffamily, anchor=west}]
  % layers from top (1) to bottom (9)
  \def\layername#1{%
    \ifcase#1\or \iflangja{アプリケーション}{applications}%
    \or \iflangja{データ}{data}%
    \or \iflangja{ランタイム}{runtime}%
    \or \iflangja{ミドルウェア}{middleware}%
    \or OS%
    \or \iflangja{仮想化}{virtualization}%
    \or \iflangja{サーバ}{servers}%
    \or \iflangja{ストレージ}{storage}%
    \or \iflangja{ネットワーク}{networking}\fi}
  % column / header / example / first layer managed by the provider (10 = none)
  \foreach \c/\h/\e/\p in {%
      0/{\iflangja{オンプレミス}{on premises}}/{}/10,
      1/IaaS/{Amazon EC2}/6,
      2/PaaS/{Google App Engine}/3,
      3/SaaS/{Gmail}/1} {
    \pgfmathsetmacro{\cx}{14 + 30*\c}
    \node[hd] at (\cx,52.3) {\h};
    \node[ex] at (\cx,46.8) {\strut\e};
    \foreach \l in {1,...,9} {
      \pgfmathsetmacro{\ly}{46 - 4.6*\l + 2.3}
      \ifnum\l<\p
        \node[own] at (\cx,\ly) {\layername{\l}};
      \else
        \node[prov] at (\cx,\ly) {\layername{\l}};
      \fi
    }
  }
  % legend
  \node[own, minimum width=5mm, minimum height=3.5mm] at (27,-3.5) {};
  \node[lab] at (30,-3.5) {\iflangja{利用者が管理}{managed by the customer}};
  \node[prov, minimum width=5mm, minimum height=3.5mm] at (70,-3.5) {};
  \node[lab] at (73,-3.5) {\iflangja{プロバイダが管理}{managed by the provider}};
\end{tikzpicture}
   (width ~116 mm)
\begin{tikzpicture}[x=1mm, y=1mm, font=\scriptsize\sffamily,
  own/.style={draw, fill=white, minimum width=26mm, minimum height=4.6mm, inner sep=0.5pt},
  prov/.style={draw, fill=ln-accent!25, minimum width=26mm, minimum height=4.6mm, inner sep=0.5pt},
  hd/.style={font=\footnotesize\sffamily\bfseries, anchor=base},
  ex/.style={font=\scriptsize\sffamily, text=ln-muted, anchor=south},
  lab/.style={font=\scriptsize\sffamily, anchor=west}]
  % layers from top (1) to bottom (9)
  \def\layername#1{%
    \ifcase#1\or \iflangja{アプリケーション}{applications}%
    \or \iflangja{データ}{data}%
    \or \iflangja{ランタイム}{runtime}%
    \or \iflangja{ミドルウェア}{middleware}%
    \or OS%
    \or \iflangja{仮想化}{virtualization}%
    \or \iflangja{サーバ}{servers}%
    \or \iflangja{ストレージ}{storage}%
    \or \iflangja{ネットワーク}{networking}\fi}
  % column / header / example / first layer managed by the provider (10 = none)
  \foreach \c/\h/\e/\p in {%
      0/{\iflangja{オンプレミス}{on premises}}/{}/10,
      1/IaaS/{Amazon EC2}/6,
      2/PaaS/{Google App Engine}/3,
      3/SaaS/{Gmail}/1} {
    \pgfmathsetmacro{\cx}{14 + 30*\c}
    \node[hd] at (\cx,52.3) {\h};
    \node[ex] at (\cx,46.8) {\strut\e};
    \foreach \l in {1,...,9} {
      \pgfmathsetmacro{\ly}{46 - 4.6*\l + 2.3}
      \ifnum\l<\p
        \node[own] at (\cx,\ly) {\layername{\l}};
      \else
        \node[prov] at (\cx,\ly) {\layername{\l}};
      \fi
    }
  }
  % legend
  \node[own, minimum width=5mm, minimum height=3.5mm] at (27,-3.5) {};
  \node[lab] at (30,-3.5) {\iflangja{利用者が管理}{managed by the customer}};
  \node[prov, minimum width=5mm, minimum height=3.5mm] at (70,-3.5) {};
  \node[lab] at (73,-3.5) {\iflangja{プロバイダが管理}{managed by the provider}};
\end{tikzpicture}

  \caption{クラウドのサービスモデルと，それぞれの例．網かけの層はプロバイダ
    が，それ以外の層は顧客が管理する．}
  \label{fig:l16-service-models}
\end{figure}

\term{IaaS}[infrastructure as a service]%
\index{infrastructure as a service|see{IaaS}}では，プロバイダが物理的な
インフラストラクチャ，すなわちネットワーク，ストレージ，サーバと，その仮想化
を管理する．顧客は仮想マシンを借り，オペレーティングシステムとそれより上の
すべてを自分で設置し管理する．Amazon EC2 がその例である．

\term{PaaS}[platform as a service]\index{platform as a service|see{PaaS}}%
では，プロバイダがオペレーティングシステム，ミドルウェア，ランタイム環境も
管理する．顧客は，プラットフォームのインタフェースに対して書かれた
アプリケーションとそのデータだけを与え，プラットフォームがアプリケーションを
配備して実行する．Google App Engine がその例である．

\term{SaaS}[software as a service]\index{software as a service|see{SaaS}}%
では，プロバイダがアプリケーションを含むスタック全体を管理し，顧客は Gmail
の利用者のようにアプリケーションを使うだけである．

\begin{example}
  \cref{sec:l16-services} の三層のオンラインショップの運用者は，仮想マシンを
  借り，そこにオペレーティングシステム，Web サーバ，アプリケーションサーバ，
  データベースサーバを設置し，これらのソフトウェアとその更新およびスケー
  リングをすべて管理することで，IaaS 上でショップを運用できる．PaaS 上では，
  運用者はビジネスロジックのコードと静的コンテンツだけをアップロードする．
  プラットフォームがそれらを実行し，管理されたデータベースを提供し，負荷が
  増えればインスタンスを追加する．同時にショップの従業員は，SaaS として提供
  される電子メールを読むかもしれない．そこで彼らが管理するのは，自分の
  アカウントとメッセージだけである．
\end{example}

\section{クラウドへの要件}
\label{sec:l16-requirements}

プロバイダの視点から見ると，クラウドコンピューティングは基盤となる技術に
いくつもの要件を課す．\source{P4L4，クラウドへの要件，クラウドにおける障害の
確率；OSTEP ch.~48．}
\begin{enumerate}
  \item \textbf{代替可能な資源．}資源は互いに交換可能であり，異なる
    オペレーティングシステムやソフトウェアスタックのように異なる要件を持つ
    顧客のために，容易に転用できなければならない．代替可能性がなければ，
    すべての顧客が専用の資源を必要とし，共有に立脚するクラウドの経済的な
    機会（\cref{sec:l16-cloud}）は存在しないことになる．
  \item \textbf{伸縮的で動的な資源割り当て．}機構と方針は，需要に応じて資源を
    割り当て，解放されたときにはそれを回収しなければならない．
  \item \textbf{規模．}管理と資源割り当ては，数千台規模のマシンからなる
    データセンタと，多数の顧客に対して機能しなければならない．
  \item \textbf{障害への対処．}この規模では障害は通常の事象であり，システムは
    それにもかかわらず動作を続けなければならない．
  \item \textbf{マルチテナンシ．}テナントは互いに隔離されなければならず，各
    テナントの性能が保証されなければならない．
  \item \textbf{セキュリティ．}各テナントのデータと計算は，他のテナントから
    も，クラウドの外部の攻撃者からも保護されなければならない．
\end{enumerate}

\begin{definition}[マルチテナンシ]
  複数の顧客が同じ物理資源を共有し，各顧客が自分の隔離された資源の集合を使う
  ことを，\term[まるちてなんし]{マルチテナンシ}[multi-tenancy]と呼ぶ．
\end{definition}

マルチテナンシは，VMM が VM に対して提供するのと同様のテナントの状態の
隔離\index{かくり@隔離}（第12章）を必要とし，さらに性能の隔離も必要とする．

\begin{definition}[性能隔離]
  あるテナントが得る性能が，資源を共有する他のテナントのワークロードに依存
  しないとき，そのシステムは\term[せいのうかくり]{性能隔離}%
  [performance isolation]を提供するという．
\end{definition}

性能隔離がなければ，ディスクやメモリを飽和させる VM と同じマシンを共有する
テナントの VM は予測できない応答時間を観測し，プロバイダは約束した
SLA\index{SLA}（第6章）を守れない．

\subsection{規模がもたらす障害}

障害は，構成要素の数が多いというだけの理由で頻繁に起こる．\caution{%
\Cref{eq:l16-failure} は構成要素が独立に故障することを仮定している．相関の
ある故障---電源や冷却の喪失，ネットワーク分断，欠陥のあるソフトウェア更新---は
多数のマシンを一度に停止させ，実際に大規模障害を引き起こすのはそれらである．}%
$N$ 個の構成要素の各々が，ある期間内に独立に確率 $p$ で故障するとすれば，その
うち少なくとも1つが故障する確率は
\begin{equation}
  P_{\mathrm{fail}} \;=\; 1 - (1-p)^{N} ,
  \label{eq:l16-failure}
\end{equation}
であり，$N$ が大きくなるにつれて1に近づく．一方，故障する構成要素の期待個数
$Np$ は $N$ に比例して増える．したがって，単体では
めったに故障しない構成要素も，データセンタのどこかでは常に故障していること
になる．クラウドのシステムは，たとえばタイムアウトによって障害を検出し，それ
にもかかわらず，これまでの章で扱った機構によって動作を続けなければならない．
故障した複製を引き継ぐ複製（第14章），計算を再開できるチェックポイント
（第8章）\index{ちぇっくぽいんてぃんぐ@チェックポインティング}，故障しつつ
あるマシンから退避させる VM のマイグレーション（第12章）である．

\begin{example}
  あるマシンが，他のマシンとは独立に，ある1日のうちに確率 $p = 0.001$ で故障
  するとする．単体では平均しておよそ3年に1度故障する計算である．
  \cref{eq:l16-failure} により，ある1日に少なくとも1台のマシンが故障する確率
  は，10台では \SI{1.0}{\percent}，100台では \SI{9.5}{\percent}，\num{1000}
  台では \SI{63}{\percent} である．\num{10000} 台では事実上
  \SI{100}{\percent} であり，平均して毎日10台が故障する．
\end{example}

\begin{keypoint}
  クラウドは，代替可能な資源，伸縮的な割り当て，規模に耐える管理，障害への
  耐性，そしてテナントの状態と性能の隔離を必要とする．構成要素が多ければ，
  ほとんど常にどこかの構成要素が故障している（\cref{eq:l16-failure}）．
\end{keypoint}

\section{クラウドを可能にする技術}
\label{sec:l16-technologies}

いくつもの技術が合わさってクラウドコンピューティングを可能にしている．
\source{P4L4，クラウドを可能にする技術；Hindman ら \cite{hindman2011}；
Dean と Ghemawat \cite{dean2004}；Zaharia ら \cite{zaharia2012}；McKeown ら
\cite{mckeown2008}．}

\paragraph{仮想化．}仮想化（第12章）は資源を代替可能にする．1台の物理マシン
は，どの顧客の VM でも，どのようなオペレーティングシステムやソフトウェア
スタックのものでも載せることができ，VM はテナントを互いに隔離する．VM は需要
の変化に応じて生成，破棄，マイグレートできる．

\paragraph{資源の割り当て．}Mesos \cite{hindman2011} や Hadoop YARN のような
クラスタの資源マネージャは，多数のマシンの資源を，その上で動作するアプリケー
ションやフレームワークの間でスケジュールし，各々がどのマシンのどの CPU と
どれだけのメモリを受け取るかを決める．\margin{Hadoop は，Google が記述した
スタックのオープンソース版に当たる．HDFS は Google File System に，Hadoop
MapReduce は MapReduce に，HBase は Bigtable に対応する．}

\paragraph{ビッグデータ処理．}Hadoop MapReduce や Spark のようなフレーム
ワークは，大規模なデータ集合を多数のマシン上で並列に処理する．

\begin{definition}[MapReduce]
  \term{MapReduce}[MapReduce] \cite{dean2004} とは，クラスタ上で大規模な
  データ集合を処理するためのプログラミングモデルであり，またそれを実装する
  ランタイムシステムである．プログラマは，各入力レコードを中間的なキーと値の
  組へ変換する \emph{map} 関数と，同じキーを持つすべての中間の値を結合する
  \emph{reduce} 関数を書く．ランタイムは入力を分割し，map タスクと reduce
  タスクを多数のマシン上で並列に実行し，中間の組をキーごとにまとめ，故障した
  マシンのタスクを再実行する．
\end{definition}

たとえば，文書の集まりの中で各単語が何回現れるかを数えるには，map 関数が単語
$w$ の出現ごとに組 $(w, 1)$ を出力し，reduce 関数が単語ごとに値を合計する．
MapReduce は中間結果をディスクに格納するが，Spark \cite{zaharia2012} はそれを
メモリに保持するので，同じデータを再利用する反復的なアルゴリズムや対話的な
問い合わせが高速になる．

\paragraph{ストレージ．}クラウドはデータを，上書きよりも追記が中心の非常に
大きなファイルのために設計された分散ファイルシステム（第14章）に，関係データ
モデルと従来のデータベースの一貫性保証の一部を手放すことでスケールアウトする
NoSQL データベースに，そして多数のマシンのメモリにデータを保持する分散
インメモリストアやキャッシュに格納する．

\paragraph{ソフトウェア定義のネットワーク，ストレージ，データセンタ．}
ネットワーク，ストレージ，そしてやがてはデータセンタ全体が，手作業ではなく
API を通じてソフトウェアによって構成される．

\begin{definition}[ソフトウェア定義ネットワーキング]
  \term{SDN}[software-defined networking]%
  \index{software-defined networking|see{SDN}}では，トラフィックをどう転送
  するかを決めるネットワークの\emph{制御プレーン}が，パケットを転送する
  スイッチの\emph{データプレーン}から分離される．ソフトウェアで実装され論理的
  に集中化されたコントローラが，OpenFlow \cite{mckeown2008} のような開かれた
  インタフェースを通じてスイッチの転送表をプログラムする．
\end{definition}

SDN によって，クラウドはネットワークをソフトウェアで再構成できる．たとえば，
各テナントに，他のテナントのものから隔離された専用の仮想ネットワークを与え
られる．\margin{OpenFlow は2008年に仕様が定められ，この種のインタフェースと
して最初に広く実装されたものであり，コントローラがスイッチの転送表にマッチ・
アンド・アクションの規則を設定できるようにする．}ソフトウェア定義ストレージと
ソフトウェア定義データセンタは，同じ原理をストレージと，データセンタのすべて
の資源に適用する．

\paragraph{監視．}数千台のマシンと，その上で動作するサービスは，ログと測定値
を絶えず生み出す．Flume，CloudWatch，Log Insight のようなリアルタイムのログ
処理システムは，障害，異常，過負荷を検出し，伸縮的なスケーリングを起動する
ために，それらを生成されるそばから収集し分析する．

\section{ビッグデータエンジンとしてのクラウド}
\label{sec:l16-bigdata}

大規模なインターネットサービスの利用者活動のログ，索引付けすべき Web ページ，
科学的な測定値など，単一のマシンには収まらないデータ集合を分析する
アプリケーションは多い．\source{P4L4，ビッグデータエンジンとしてのクラウド，
ビッグデータスタックの例．}

\begin{definition}[ビッグデータ]
  単一のマシンが格納し処理できる大きさを超え，多数のマシン上で格納し分析
  しなければならないデータ集合を，\term[びっぐでーた]{ビッグデータ}%
  [big data]と呼ぶ．
\end{definition}

クラウドは，そうしたデータのストレージと，それを並列に処理する多数のマシンの
両方を提供する．それを行うソフトウェアは，次の層からなる\emph{ビッグデータ
エンジン}を成す（\cref{fig:l16-bigdata}）．
\begin{description}
  \item[データストレージ] データを多数のマシンに分散し複製して格納する．
  \item[データ処理] データに対する計算を多数のマシン上で並列に実行する．
  \item[キャッシュ] 頻繁に使うデータをメモリに保持し，ストレージ層の
    レイテンシを避ける．
  \item[言語フロントエンド] 利用者が処理層を直接プログラムする代わりに，
    SQL 風の問い合わせ言語のような高水準の言語でデータを問い合わせ分析できる
    ようにする．
  \item[分析ライブラリ] 機械学習，グラフ分析，統計などのために，処理層の上に
    構築されたアルゴリズムを提供する．
  \item[ストリームデータ] ログ，クリック，センサの測定値といったデータは
    絶えずエンジンへ入ってくるので，格納された後だけでなく，到着するそばから
    処理されなければならない．
\end{description}

\begin{figure}[htbp]
  \centering
  % The cloud as a big data engine: the layers of a big data stack, and the
% streaming data that enter it continuously.
% Usage: % The cloud as a big data engine: the layers of a big data stack, and the
% streaming data that enter it continuously.
% Usage: % The cloud as a big data engine: the layers of a big data stack, and the
% streaming data that enter it continuously.
% Usage: \input{figures/l16-bigdata}   (width ~116 mm)
\begin{tikzpicture}[x=1mm, y=1mm, font=\scriptsize\sffamily,
  >={Stealth[length=1.5mm]},
  ly/.style={draw, fill=ln-box, minimum height=8mm, inner sep=1pt, align=center},
  top/.style={draw, fill=ln-accent!18, minimum height=8mm, inner sep=1pt, align=center},
  st/.style={draw, fill=ln-accent!35, minimum height=8mm, inner sep=1pt, align=center},
  lab/.style={font=\scriptsize\sffamily, text=ln-muted, align=center}]
  % users of the engine
  \node[lab] at (74,49) {\iflangja{アプリケーション，データ分析者}{applications, data analysts}};
  \draw[<->, thick] (57,46.5) -- (57,42.5);
  \draw[<->, thick] (95,46.5) -- (95,42.5);
  % layers
  \node[top, minimum width=42mm] (lang) at (53,38) {\iflangja{言語フロントエンド}{language front-ends}\\\textcolor{ln-muted}{\iflangja{例: 問い合わせ}{e.g.\ queries}}};
  \node[top, minimum width=42mm] (lib) at (97,38) {\iflangja{分析ライブラリ}{analytics libraries}\\\textcolor{ln-muted}{\iflangja{例: 機械学習}{e.g.\ machine learning}}};
  \node[ly, minimum width=86mm] (proc) at (75,27) {\iflangja{データ処理}{data processing}\\\textcolor{ln-muted}{\iflangja{多数のマシンで並列に実行}{in parallel on many machines}}};
  \node[ly, minimum width=86mm] (cache) at (75,16) {\iflangja{キャッシュ}{caching}\\\textcolor{ln-muted}{\iflangja{頻繁に使うデータをメモリに保持}{frequently used data kept in memory}}};
  \node[st, minimum width=86mm] (store) at (75,5) {\iflangja{データストレージ}{data storage}\\\textcolor{ln-muted}{\iflangja{多数のマシンに分散して格納}{distributed over many machines}}};
  % streaming data
  \node[lab, text=ln-accent, anchor=east] (sd) at (27,27)
    {\iflangja{連続して到着する}{continuously}\\\iflangja{ストリームデータ}{streaming data}};
  \draw[->, line width=1.2pt, ln-accent] (sd.east) -- (proc.west);
  \draw[->, line width=1.2pt, ln-accent] (29.2,27) |- (store.west);
  \fill[ln-accent] (29.2,27) circle (0.8);
\end{tikzpicture}
   (width ~116 mm)
\begin{tikzpicture}[x=1mm, y=1mm, font=\scriptsize\sffamily,
  >={Stealth[length=1.5mm]},
  ly/.style={draw, fill=ln-box, minimum height=8mm, inner sep=1pt, align=center},
  top/.style={draw, fill=ln-accent!18, minimum height=8mm, inner sep=1pt, align=center},
  st/.style={draw, fill=ln-accent!35, minimum height=8mm, inner sep=1pt, align=center},
  lab/.style={font=\scriptsize\sffamily, text=ln-muted, align=center}]
  % users of the engine
  \node[lab] at (74,49) {\iflangja{アプリケーション，データ分析者}{applications, data analysts}};
  \draw[<->, thick] (57,46.5) -- (57,42.5);
  \draw[<->, thick] (95,46.5) -- (95,42.5);
  % layers
  \node[top, minimum width=42mm] (lang) at (53,38) {\iflangja{言語フロントエンド}{language front-ends}\\\textcolor{ln-muted}{\iflangja{例: 問い合わせ}{e.g.\ queries}}};
  \node[top, minimum width=42mm] (lib) at (97,38) {\iflangja{分析ライブラリ}{analytics libraries}\\\textcolor{ln-muted}{\iflangja{例: 機械学習}{e.g.\ machine learning}}};
  \node[ly, minimum width=86mm] (proc) at (75,27) {\iflangja{データ処理}{data processing}\\\textcolor{ln-muted}{\iflangja{多数のマシンで並列に実行}{in parallel on many machines}}};
  \node[ly, minimum width=86mm] (cache) at (75,16) {\iflangja{キャッシュ}{caching}\\\textcolor{ln-muted}{\iflangja{頻繁に使うデータをメモリに保持}{frequently used data kept in memory}}};
  \node[st, minimum width=86mm] (store) at (75,5) {\iflangja{データストレージ}{data storage}\\\textcolor{ln-muted}{\iflangja{多数のマシンに分散して格納}{distributed over many machines}}};
  % streaming data
  \node[lab, text=ln-accent, anchor=east] (sd) at (27,27)
    {\iflangja{連続して到着する}{continuously}\\\iflangja{ストリームデータ}{streaming data}};
  \draw[->, line width=1.2pt, ln-accent] (sd.east) -- (proc.west);
  \draw[->, line width=1.2pt, ln-accent] (29.2,27) |- (store.west);
  \fill[ln-accent] (29.2,27) circle (0.8);
\end{tikzpicture}
   (width ~116 mm)
\begin{tikzpicture}[x=1mm, y=1mm, font=\scriptsize\sffamily,
  >={Stealth[length=1.5mm]},
  ly/.style={draw, fill=ln-box, minimum height=8mm, inner sep=1pt, align=center},
  top/.style={draw, fill=ln-accent!18, minimum height=8mm, inner sep=1pt, align=center},
  st/.style={draw, fill=ln-accent!35, minimum height=8mm, inner sep=1pt, align=center},
  lab/.style={font=\scriptsize\sffamily, text=ln-muted, align=center}]
  % users of the engine
  \node[lab] at (74,49) {\iflangja{アプリケーション，データ分析者}{applications, data analysts}};
  \draw[<->, thick] (57,46.5) -- (57,42.5);
  \draw[<->, thick] (95,46.5) -- (95,42.5);
  % layers
  \node[top, minimum width=42mm] (lang) at (53,38) {\iflangja{言語フロントエンド}{language front-ends}\\\textcolor{ln-muted}{\iflangja{例: 問い合わせ}{e.g.\ queries}}};
  \node[top, minimum width=42mm] (lib) at (97,38) {\iflangja{分析ライブラリ}{analytics libraries}\\\textcolor{ln-muted}{\iflangja{例: 機械学習}{e.g.\ machine learning}}};
  \node[ly, minimum width=86mm] (proc) at (75,27) {\iflangja{データ処理}{data processing}\\\textcolor{ln-muted}{\iflangja{多数のマシンで並列に実行}{in parallel on many machines}}};
  \node[ly, minimum width=86mm] (cache) at (75,16) {\iflangja{キャッシュ}{caching}\\\textcolor{ln-muted}{\iflangja{頻繁に使うデータをメモリに保持}{frequently used data kept in memory}}};
  \node[st, minimum width=86mm] (store) at (75,5) {\iflangja{データストレージ}{data storage}\\\textcolor{ln-muted}{\iflangja{多数のマシンに分散して格納}{distributed over many machines}}};
  % streaming data
  \node[lab, text=ln-accent, anchor=east] (sd) at (27,27)
    {\iflangja{連続して到着する}{continuously}\\\iflangja{ストリームデータ}{streaming data}};
  \draw[->, line width=1.2pt, ln-accent] (sd.east) -- (proc.west);
  \draw[->, line width=1.2pt, ln-accent] (29.2,27) |- (store.west);
  \fill[ln-accent] (29.2,27) circle (0.8);
\end{tikzpicture}

  \caption{ビッグデータエンジンの層．ストリームデータは絶えずエンジンへ入り，
    到着するそばから処理され，かつ格納される．}
  \label{fig:l16-bigdata}
\end{figure}

広く使われている2つのオープンソースのスタックが，これらの層を実装している
（\cref{tab:l16-stacks}）．\emph{Hadoop} スタックは，Hadoop Distributed File
System（HDFS）と MapReduce を中心に構築されている．カリフォルニア大学
バークレー校で開発された \emph{Berkeley Data Analytics Stack}（BDAS）は，
Spark とクラスタの資源マネージャ Mesos を中心に構築されており，その下で
HDFS や他のストレージシステムを利用できる．\margin{いくつかの構成要素はその後
に改称あるいは置き換えられている．Tachyon は2016年に Alluxio となり，Shark は
Spark SQL に取って代わられた．\Cref{tab:l16-stacks} は元の名称のままである．}

\begin{table}[htbp]
  \centering
  \caption{2つのビッグデータスタックの構成要素（層ごと）．}
  \label{tab:l16-stacks}
  \small
  \begin{tabular}{@{}>{\raggedright\arraybackslash}p{0.2\linewidth}>{\raggedright\arraybackslash}p{0.36\linewidth}>{\raggedright\arraybackslash}p{0.36\linewidth}@{}}
    \toprule
    層 & Hadoop & BDAS \\
    \midrule
    データストレージ & HDFS（ファイルシステム），HBase（列指向ストア） &
      HDFS，Amazon S3 その他のストレージシステム \\
    キャッシュ & --- & Tachyon（インメモリストレージ） \\
    資源管理 & YARN & Mesos，YARN \\
    データ処理 & MapReduce & Spark \\
    ストリームデータ & Flume（ログ収集） & Spark Streaming \\
    言語フロントエンド & Hive（SQL），Pig（データフロー） & Shark（SQL），
      BlinkDB（近似問い合わせ） \\
    分析ライブラリ & Mahout（機械学習），R コネクタ（統計） & MLlib，MLbase
      （機械学習），GraphX（グラフ） \\
    協調 & ZooKeeper（協調），Oozie（ワークフロー），Sqoop（関係データベース
      とのデータ交換） & --- \\
    \bottomrule
  \end{tabular}
\end{table}

\needspace{8\baselineskip}
\begin{keypoint}[章のまとめ]
  インターネットサービスは，プレゼンテーション，ビジネスロジック，データ
  ベースの層からなり，ロードバランサの背後で多数のノードへスケールアウト
  される．ホモジニアスなノードは管理が単純であり，ヘテロジニアスなノードは
  局所性を活かせるが，要求を調べるフロントエンドと需要に従う管理を要する．
  クラウドコンピューティングは，共有された資源を API を通じて伸縮的に，利用に
  応じた課金で提供し，プロバイダが管理する．これが成り立つのは大数の法則と
  規模の経済のゆえであり，ユーティリティコンピューティングの構想に近づく．
  クラウドはパブリック，プライベート，ハイブリッド，コミュニティのいずれかで
  あり，インフラストラクチャ，プラットフォーム，ソフトウェアをサービスとして
  提供する．クラウドが要するのは，代替可能な資源，規模に耐える伸縮的な割り
  当て，頻繁な障害への耐性，テナントの隔離であり，それらを与えるのが仮想化，
  クラスタの資源管理，MapReduce のようなビッグデータフレームワーク，分散
  ストレージ，ソフトウェア定義ネットワーキング，監視である．これらの技術の
  上で，クラウドはビッグデータエンジンとして働く．
\end{keypoint}

\end{document}
